\chapter{Anatomy of a Stat-Arb Book}\label{ch:s1:anatomy-of-a-stat-arb-book}

On an average day a book holds 452 stocks long and 495 short. Its median position is a quarter of one per cent of its capital, and no position is larger than one per cent.
Over eight simulated years its best day makes 1.5 per cent and its worst loses 1.25, and its profit is the sum of 1.9 million position-days, none of which
matters much alone. Remove one line of its construction (the neutrality to the style factors) and the same signals give a book whose best and worst days are
nearly four times larger and whose daily profit and loss is, to 99.7 per cent of its variance, a bet on one factor. This part of the book is about strategies; its
first chapter is about the vehicle most of the equity strategies ride in, the market-neutral \omterm{def:s1:anatomy-of-a-stat-arb-book:statarb}{stat-arb book}: how it is built from Book~7's pieces, what makes it
neutral, how it is financed, and where each day's P\&L comes from. The build is \texttt{firm.statbook}.

\section{Thousands of small bets}

\begin{definition}[Statistical arbitrage, stat-arb book]\label{def:s1:anatomy-of-a-stat-arb-book:statarb}
\emph{Statistical arbitrage}\index{statistical arbitrage} is the systematic trading of many small, individually unreliable relative mispricings among
securities, each expected to earn a little, so that the book's return comes from their number. A \emph{stat-arb book}\index{stat-arb book} is a long--short
portfolio of hundreds or thousands of stocks built this way, neutral to the market and usually to industries and styles, rebuilt daily or more often.
\end{definition}

The fundamental law of active management (Book~7, chapter~15) is the \omterm{def:s1:anatomy-of-a-stat-arb-book:statarb}{stat-arb book}'s business plan: an information ratio grows with the square root of
breadth, so a skill too small to trade in one stock becomes a business across a thousand, rebuilt every day. The price is everything the law assumes away:
the bets must be independent, which a book's common exposures prevent, and cheap, which daily trading in a thousand names is not.

The chapter's book lives on \texttt{firm.synthmkt} (Book~7, chapter~5), all of its listed names, years~3 to~10. Its signal is an equal blend of the two
predictors Book~7 measured best: the one-day residual reversal and 12--1 month momentum, both point in time, smoothed with a half-life of two days to slow
its trading. Each close, the signal is made neutral by regression on the risk model's exposures, weighted to a gross exposure of three times capital (1.5
long, 1.5 short), capped at 1\% of capital a name, and held over the next day at a cost of five basis points of the weight traded (\cref{lst:s1:anatomy-of-a-stat-arb-book:book}).

\section{Neutrality: dollar, beta, sector, factor}

Neutrality is a set of constraints on the book's exposures (dollar, beta and factor neutrality are Book~7, chapter~25's terms). Their purpose is that the book's
P\&L should come from the stock-specific part of its signal and not from the market or a style. The chapter builds two versions of the same book. The first is
neutral to the dollar, the market's beta and the ten industries, which is what many descriptions of market-neutral books promise; the second also to the size,
value and momentum styles.

\begin{center}
\small
\begin{tabular}{@{}lrr@{}}
\toprule
years 3 to 10, \$1 billion of capital & neutral to beta and industries & also to the styles\\
\midrule
return a year, after costs and financing & 23.6\% & 22.5\%\\
volatility & 23.6\% & 4.5\%\\
Sharpe ratio & 1.00 & 5.03\\
share of the P\&L variance from factors & 99.7\% & 7.6\%\\
best and worst day & $+5.8\%$, $-4.8\%$ & $+1.5\%$, $-1.25\%$\\
one-way turnover a day & 42\% & 91\%\\
trading costs a year & $-10.6\%$ & $-22.9\%$\\
\bottomrule
\end{tabular}
\end{center}

The first book is neutral to what it was asked to be neutral to and not to what it holds. Its momentum signal gives it a large exposure to the momentum style,
and the momentum factor's contribution alone has a volatility of 23.1\% a year: the book's specific P\&L is swamped, and 99.7\% of its variance is factor
risk (\cref{fig:s1:anatomy-of-a-stat-arb-book:cum}). Neutralised to the styles as well, the book keeps its specific return and loses the factor bet: its
volatility falls to 4.5\% and the factors explain 7.6\% of its variance. The Sharpe ratio of 5 belongs to the simulation, which plants a strong reversal and a
persistent drift (Book~7, chapter~5): a real book would earn a fraction of it, as the public record in section~4 shows. What carries over is the shape of the
comparison. The neutral book pays for its neutrality in turnover, 91\% of capital a day against 42\%, because the style exposures move every day and must be
traded away.

\begin{omfigure}
\begin{tikzpicture}
\begin{axis}[name=a,width=6.8cm,height=5.4cm,xlabel={\small year},ylabel={\small cumulative P\&L (fraction of capital)},tick label style={font=\footnotesize},
  xmin=2,xmax=10,ymin=-1,ymax=4,grid=major,grid style={gray!25},title={\small neutral to beta and industries},legend pos=north west,legend style={font=\scriptsize},
  table/col sep=comma]
\addplot[omDef,very thick] table[x=year,y=partial_specific]{figdata/strategies-1/01-anatomy-of-a-stat-arb-book/cum.csv};
\addplot[omThm,thick] table[x=year,y=partial_factor]{figdata/strategies-1/01-anatomy-of-a-stat-arb-book/cum.csv};
\addplot[black!60,thick,dashed] table[x=year,y=partial_total]{figdata/strategies-1/01-anatomy-of-a-stat-arb-book/cum.csv};
\legend{specific,factor,net of costs and financing}
\end{axis}
\begin{axis}[at={(a.east)},anchor=west,xshift=1.2cm,width=6.8cm,height=5.4cm,xlabel={\small year},tick label style={font=\footnotesize},
  xmin=2,xmax=10,ymin=-1,ymax=4,grid=major,grid style={gray!25},title={\small also neutral to the styles},table/col sep=comma]
\addplot[omDef,very thick] table[x=year,y=full_specific]{figdata/strategies-1/01-anatomy-of-a-stat-arb-book/cum.csv};
\addplot[omThm,thick] table[x=year,y=full_factor]{figdata/strategies-1/01-anatomy-of-a-stat-arb-book/cum.csv};
\addplot[black!60,thick,dashed] table[x=year,y=full_total]{figdata/strategies-1/01-anatomy-of-a-stat-arb-book/cum.csv};
\end{axis}
\end{tikzpicture}

\omcaption{Cumulative daily P\&L of the two books, attributed to the risk model's factors and to specific returns, and net of trading costs and financing.
Left, the momentum factor's swings dominate; right, the factor line stays near zero. Data: \texttt{s1\_statbook.run}.}\label{fig:s1:anatomy-of-a-stat-arb-book:cum}
\end{omfigure}

\section{Financing: cash, leverage, borrow and rebates}

\begin{definition}[Hard-to-borrow list]\label{def:s1:anatomy-of-a-stat-arb-book:htb}
A prime broker's \emph{hard-to-borrow list}\index{hard-to-borrow list} names the stocks it cannot lend at the general-collateral fee: shorting them costs a
higher borrow fee (Book~1, chapter~6), may be refused, and is exposed to recall.
\end{definition}

A long--short book is a financing arrangement as much as a portfolio. The longs are paid for with the fund's capital and, beyond it, with cash borrowed from the
prime broker; the shorts are sold with borrowed stock, and the sale proceeds stay with the broker as collateral, earning the rebate rate (Book~1, chapter~6),
the benchmark rate less the stock's borrow fee. \texttt{firm.statbook.financing} (\cref{lst:s1:anatomy-of-a-stat-arb-book:fin}) charges one day of each: a
debit balance of 0.5 of capital (longs of 1.5 on capital of 1) at the benchmark rate of 4\% plus a 0.5\% spread, and short proceeds of 1.5 at 4\% less the fee,
with 95\% of names at a general-collateral fee of 0.25\% and 5\%, the \omterm{def:s1:anatomy-of-a-stat-arb-book:htb}{hard-to-borrow list}, at fees drawn around 3\%, not shorted by the book. For the neutral
book the rebate brings 5.6\% of capital a year and the long financing costs 2.2\%; the borrow fees inside the rebate cost 0.4\%. The net financing is $+3.4\%$,
less than the 4\% the capital would earn in cash, and the long financing and borrow fees together are 6.2\% of the book's gross alpha of 42.0\%.

Gross exposure of three times capital is not allowed under the strategy-based margin rules of retail accounts; it is a professional book's leverage, reached with
risk-based margin (\cref{dat:s1:anatomy-of-a-stat-arb-book:margin}).

\begin{dated}{2026-09}{Margin for a long--short equity book}\label{dat:s1:anatomy-of-a-stat-arb-book:margin}
In the United States, Regulation~T (12 CFR 220.12) requires margin of 50\% of a long equity position's value and 150\% of a short sale's value, of which the
proceeds supply 100\%: a book's equity must be half its gross, so gross exposure is at most twice capital; the chapter's book, at three times, would need equity of
1.5 times its capital. FINRA Rule 4210 lets members use portfolio margin instead, in which requirements come from a stress of each
position (plus and minus 15\% for an equity security) rather than fixed percentages; a book stressed at 15\% on every position without netting needs 0.45 of
capital for a gross of three. Strategy-based maintenance on a short stock priced at \$5 or more is the greater of \$5 a share and 30\% of its value.
\end{dated}

\section{What public descriptions of such books say}

The published record of stat-arb performance is thin and dated, and it says two things. Avellaneda and Lee, trading residuals against principal components and
sector funds across the US market, reported Sharpe ratios after costs of 1.44 for 1997--2007, stronger before 2003 and 0.9 in 2003--2007, with the ETF version
degrading similarly after 2002. Khandani and Lo documented the other side of the business: many books built this way held similar positions, and in August 2007
they lost together (Book~7, chapter~28). Pedersen's book on hedge-fund strategies is the practitioner's account to read alongside them. The chapters that follow take its ingredients one at a time: reversal
(chapter~2), residuals (chapter~3), pairs (chapter~4), momentum (chapter~5) and the rest of Part~I.

\section{The daily cycle of a stat-arb book}

\begin{definition}[Daily P\&L attribution]\label{def:s1:anatomy-of-a-stat-arb-book:attribution}
\emph{Daily P\&L attribution}\index{daily P\&L attribution} splits a book's P\&L for the day into the part explained by its factor exposures times the risk
model's factor returns, the specific part (its weights times the stocks' specific returns), trading costs and financing, so that the parts add up to the P\&L.
\end{definition}

The book's day has a fixed rhythm. After the close, the day's data are checked and stored point in time; the risk model is updated; the signals are computed,
blended and smoothed; the book is optimised under its neutralities and limits, and the trade list goes to execution for the next session; overnight the financing
accrues; the next evening the day's P\&L is attributed and reconciled with the prime broker's statement. Attribution is where a desk sees that a book is not what it
thinks: with the risk model's factor returns $f_t$, exposures $X$ and specific returns $e_t$, the day's return $w^\top r_t = (X^\top w)^\top f_t + w^\top e_t$ holds
exactly, so the split is an identity, and the checks of the chapter's run confirm that the attributed parts add up to the P\&L to within $10^{-15}$ of capital.

\begin{omfigure}
\begin{tikzpicture}[st/.style={draw=omDef,fill=omDef!8,rounded corners=2pt,minimum width=1.95cm,align=center,font=\scriptsize,minimum height=1.0cm},
  arr/.style={-{Stealth},thick,omDef}]
\node[st] (d) at (0,0) {close:\\data checked};
\node[st] (r) at (2.35,0) {risk model\\updated};
\node[st] (s) at (4.7,0) {signals\\blended};
\node[st] (o) at (7.05,0) {book\\optimised};
\node[st] (e) at (9.4,0) {trades\\next day};
\node[st] (a) at (11.75,0) {P\&L attributed,\\reconciled};
\foreach \x/\y in {d/r,r/s,s/o,o/e,e/a} \draw[arr] (\x) -- (\y);
\draw[omThm,thick,-{Stealth},dashed] (a.south) to[bend left=14] node[below,font=\scriptsize,text=omThm] {financing accrues overnight; the next close starts again} (d.south);
\end{tikzpicture}

\omcaption{The daily cycle of a \omterm{def:s1:anatomy-of-a-stat-arb-book:statarb}{stat-arb book}.}\label{fig:s1:anatomy-of-a-stat-arb-book:cycle}
\end{omfigure}

\section{Tutorial: three million small bets}\label{tut:s1:anatomy-of-a-stat-arb-book:tut}

\begin{tutorial}
\textbf{Goal.} Build the two books on \texttt{firm.synthmkt}, finance them and attribute their daily P\&L. \textbf{End state:} the table and
\cref{fig:s1:anatomy-of-a-stat-arb-book:cum}.
\begin{tutsteps}
\item \textbf{The book}: blend, smooth, neutralise, cap, avoid hard-to-borrow shorts.
\omcode{code/strategies-1/01-anatomy-of-a-stat-arb-book/python/s1_statbook.py}{55}{70}{One close of the book's construction.}\label{lst:s1:anatomy-of-a-stat-arb-book:book}
\item \textbf{Financing and attribution.}
\omcode{code/firm/statbook/firm_statbook.py}{50}{74}{A day's financing, and the attribution that adds up.}\label{lst:s1:anatomy-of-a-stat-arb-book:fin}
\item \textbf{Run} \texttt{s1\_statbook.summary} for both books, \texttt{bets}, \texttt{margin\_example} and \texttt{fig\_statbook.py}.
\end{tutsteps}
\textbf{What to change next.} Raise the borrow fees of the hard-to-borrow names and let the book short them; neutralise only the momentum style and see how much of
the factor risk remains.
\end{tutorial}

\section{Build: stat-arb book accounting}\label{bld:s1:anatomy-of-a-stat-arb-book:statbook}

\begin{build}
\bfield{Purpose} The accounting under every long--short strategy of the book: financing by name, margin, and a daily attribution that reconciles to the P\&L.
\bfield{Interface} \texttt{borrow\_fees(n, seed, htb\_share, gc\_fee, htb\_median)}, \texttt{hard\_to\_borrow(fees, threshold)}, \texttt{reg\_t\_equity(long\_mv,
short\_mv)}, \texttt{stress\_equity(w, capital, move)}, \texttt{financing(w, capital, rate, long\_spread, fees)}, \texttt{attribution(w, X, f, e, cost, fin)}.
\bfield{Rules} Financing is charged by name, every day; hard-to-borrow names are flagged before the optimiser sees them; the attribution must add up to the
P\&L, or the day is not closed.
\bfield{Acceptance tests} \texttt{code/firm/statbook/tests/}: the fee distribution; Regulation~T and the stress requirement by hand; a day's financing by hand; an
attribution that adds up to $w^\top(Xf + e)$ less costs plus financing.
\bfield{Stretch} Locates and recalls day by day; margin under a portfolio-margin stress with netting; reconciliation against a broker statement.
\end{build}

\begin{omsources}
\item M.~Avellaneda and J.-H.~Lee, ``Statistical arbitrage in the US equities market'', \emph{Quantitative Finance} 10(7), 2010.
\item A.~E.~Khandani and A.~W.~Lo, ``What happened to the quants in August 2007?'', \emph{Journal of Financial Markets} 14(1), 2011.
\item L.~H.~Pedersen, \emph{Efficiently Inefficient}, Princeton University Press, 2015.
\item Regulation~T, 12 CFR 220.12; FINRA Rule 4210.
\end{omsources}

\section{Exercises}

\begin{exercise}[$\star$]\label{exo:s1:anatomy-of-a-stat-arb-book:1}
A book holds 946 positions and trades every day for a year; if each position-day were an independent bet with a standard deviation of 2\% of the position, what
would be the standard deviation, in units of one position, of the year's summed P\&L?
\end{exercise}

\begin{exercise}[$\star$]\label{exo:s1:anatomy-of-a-stat-arb-book:2}
The neutral book's net financing is 3.42\% a year. What does it cost relative to capital left in cash at the 4\% benchmark?
\end{exercise}

\begin{exercise}[$\star$]\label{exo:s1:anatomy-of-a-stat-arb-book:3}
What equity does Regulation~T require for a book long 1.5 and short 1.5 per unit of capital? What is the largest gross exposure it allows, and what a 15\% stress
without netting?
\end{exercise}

\begin{exercise}[$\star\star$]\label{exo:s1:anatomy-of-a-stat-arb-book:4}
A book is long 1.5 and short 1.5, all shorts at the general-collateral fee of 0.25\%, the benchmark at 4\% and the long spread at 0.5\%. Compute its financing for a
year.
\end{exercise}

\begin{exercise}[$\star\star$]\label{exo:s1:anatomy-of-a-stat-arb-book:5}
Why did neutralising the styles double the book's turnover?
\end{exercise}

\begin{exercise}[$\star\star$]\label{exo:s1:anatomy-of-a-stat-arb-book:6}
Show that the attribution $w^\top r = (X^\top w)^\top f + w^\top e$ is an identity when $e$ is the cross-sectional regression's residual, and explain what would break
it.
\end{exercise}

\begin{exercise}[$\star\star\star$]\label{exo:s1:anatomy-of-a-stat-arb-book:7}
\emph{Coding.} Run \texttt{s1\_statbook.summary} with a half-life of smoothing of one day and of five, and report the Sharpe ratio and the trading costs.
\end{exercise}

\begin{exercise}[$\star\star\star$]\label{exo:s1:anatomy-of-a-stat-arb-book:8}
\emph{Find the flaw.} ``Our book is market neutral: its beta is zero every day and it is dollar neutral, so its returns are pure alpha.''
\end{exercise}

\section{Problem: Three Million Small Bets}

\begin{problem}[{Weekend problem --- a book, taken apart}]\label{pb:s1:anatomy-of-a-stat-arb-book:1}
The chapter's two books on \texttt{firm.synthmkt}.

\medskip\noindent\textbf{Part I --- The book.}
\begin{enumerate}
\item What is the signal, and how is it made into weights?
\item How many names does the neutral book hold, and how large is its median position?
\item Why does breadth matter, and what does the fundamental law assume that the book violates?
\item How many position-days make its eight years?
\end{enumerate}

\medskip\noindent\textbf{Part II --- Neutrality.}
\begin{enumerate}[resume]
\item Give the two books' returns, volatilities and Sharpe ratios.
\item What share of each book's variance comes from factors, and from which factor in the first?
\item What does the neutral book pay for its neutrality?
\item Why is a Sharpe ratio of 5 a property of the simulation?
\end{enumerate}

\medskip\noindent\textbf{Part III --- Financing.}
\begin{enumerate}[resume]
\item Break the neutral book's financing into rebate, long financing and borrow fees.
\item What does Regulation~T allow, and how does a book reach a gross of three?
\item How are hard-to-borrow names handled, and what do they cost?
\end{enumerate}

\medskip\noindent\textbf{Part IV --- The verdict.}
\begin{enumerate}[resume]
\item State the \emph{named result}: the share of the year's P\&L variance explained by factor exposures in the book that was not neutral to the styles, and the
financing drag as a fraction of gross alpha.
\item What did Avellaneda and Lee find, and what happened after 2002?
\item What did August 2007 show about books built this way?
\item Describe the daily cycle.
\item Why must the attribution add up?
\item What would you add to make the neutral book realistic?
\item How would you tell a factor bet disguised as stat arb?
\item What is the first question to ask about any market-neutral fund's returns?
\item In one sentence: what is a \omterm{def:s1:anatomy-of-a-stat-arb-book:statarb}{stat-arb book}?
\end{enumerate}
\end{problem}

\section{Interview questions}

\begin{interviewq}[$\star$ \iqroles{researcher, trader}]\label{iq:s1:anatomy-of-a-stat-arb-book:1}
What does it mean for a long--short equity book to be market neutral, and what is it still exposed to?
\end{interviewq}

\begin{interviewq}[$\star\star$ \iqroles{risk}]\label{iq:s1:anatomy-of-a-stat-arb-book:2}
How is a long--short equity book financed, and what can go wrong with the financing?
\end{interviewq}

\begin{interviewq}[$\star\star$ \iqroles{researcher}]\label{iq:s1:anatomy-of-a-stat-arb-book:3}
Explain the fundamental law of active management and why it favours stat arb.
\end{interviewq}

\begin{interviewq}[$\star\star$ \iqroles{researcher, trader}]\label{iq:s1:anatomy-of-a-stat-arb-book:4}
Published stat-arb Sharpe ratios fell after the early 2000s. Why might that be?
\end{interviewq}

\begin{interviewq}[$\star\star$ \iqroles{developer, risk}]\label{iq:s1:anatomy-of-a-stat-arb-book:5}
Design the \omterm{def:s1:anatomy-of-a-stat-arb-book:attribution}{daily P\&L attribution} for a book of a thousand names. What must reconcile?
\end{interviewq}

\begin{interviewq}[$\star\star\star$ \iqroles{researcher}]\label{iq:s1:anatomy-of-a-stat-arb-book:6}
A book has gross exposure $G$ with $N$ equal positions, specific volatility $\sigma$ per name and independent specific returns. Derive its specific volatility and the
information ratio for a per-name edge $\alpha$; how does it scale with $N$?
\end{interviewq}
